\documentclass[11pt,a4paper]{article}
\usepackage[english]{babel}
\usepackage[nottoc]{tocbibind}

\usepackage{pgfplots}
\pgfplotsset{compat=1.18}
% =========================
% Packages
% =========================
\usepackage[margin=1in]{geometry}
\usepackage{amsmath, amssymb, amsthm}
\usepackage{mathtools}
\usepackage{physics}
\usepackage{bm}
\usepackage{mathrsfs}

\usepackage{enumitem}
\usepackage{hyperref}
\usepackage{xcolor}
\usepackage{fancyhdr}
\usepackage{setspace}
\usepackage{graphicx}
\usepackage{xfrac}


\usepackage{forest}
\usepackage{tikz-qtree}
\usepackage{tikz}
\usetikzlibrary{arrows.meta}
% =========================
% Page Style
% =========================
\pagestyle{fancy}
\fancyhf{}
\lhead{\textsc{Lecture Notes}}
\rhead{\textsc{MAT211}}
\cfoot{\thepage}

\setstretch{1.15}

% =========================
% Custom Colors
% =========================
\definecolor{theoremcolor}{RGB}{0,0,120}
\definecolor{definitioncolor}{RGB}{0,90,0}
\definecolor{examplecolor}{RGB}{120,0,0}

% =========================
% Theorem Environments
% =========================
\theoremstyle{definition}
\newtheorem{definition}{Definition}[section]
\newtheorem{example}[definition]{Example}
\newtheorem{remark}[definition]{Remark}
\newtheorem{problem}[definition]{Problem}

\theoremstyle{plain}
\newtheorem{theorem}[definition]{Theorem}
\newtheorem{proposition}[definition]{Proposition}
\newtheorem{lemma}[definition]{Lemma}
\newtheorem{corollary}[definition]{Corollary}

% =========================
% Custom Commands
% =========================
\newcommand{\R}{\mathbb{R}}
\newcommand{\C}{\mathbb{C}}
\newcommand{\N}{\mathbb{N}}
\newcommand{\Z}{\mathbb{Z}}
\newcommand{\Q}{\mathbb{Q}}

\newcommand{\eps}{\varepsilon}

% =========================
% Title Info
% =========================
\title{\textbf{MAT111: Principles of Mathematics}\\
\large Lecture Notes}
\author{Instructor: Emon Hossain}
\date{\today}

% =========================
% Document
% =========================
\begin{document}

\maketitle
\vspace{-1em}

\tableofcontents
\newpage

% =========================
% Lecture Starts
% =========================
\section{Lecture-01}

Aim: Provide a proper mathematical language for 
\begin{itemize}
    \item Classical Mechanics
    \item Electromagnetism
    \item Quantum Mechanics
    \item Statistical Mechanics
\end{itemize}


\subsection{Set}

\subsection{Cardinality}


The interesting number paradox is a humorous paradox which arises from the attempt to classify 
every natural number as either "interesting" or "uninteresting". The paradox states that every 
natural number is interesting.The "proof" is by contradiction: if there exists a non-empty set 
of uninteresting natural numbers, there would be a smallest uninteresting number - 
but the smallest uninteresting number is itself interesting because it is the smallest 
uninteresting number, thus producing a contradiction.

\begin{definition}
\label{D1}
A set is finite if there is a one-to-one correspondence with the set $\{1,2,3,…,n\}$ 
for some natural number $n$ ($n=0$ gives the empty set as a special case).
\end{definition}

\begin{definition}
\label{D2}
A set is finite if there is no one-to-one map of the set to a proper subset of itself 
(this is also called Dedekind finite).
\end{definition}

\begin{definition}
Hereditarily finite sets are defined as finite sets whose elements are all hereditarily finite sets. 
In other words, the set itself is finite, and all of its elements are finite sets, 
recursively all the way down to the empty set.
\end{definition}

Now, in standard set theory, these two notions are equivalent, so any set that is finite via 
\ref{D1} is finite via \ref{D2} and vice versa.

BUT, this equivalence requires the axiom of choice (well, a countable version of it). 
If you like to do set theory without that axiom, these two notions are different. 
It is possible to be finite via \ref{D2} but not finite via \ref{D1} (the reverse is impossible). 
Such a set is called a Dedekind set. So, by \ref{D2}, it is finite, but by \ref{D1} it is infinite.

Such sets can exist without the axiom of choice. For example, it is possible to have an infinite 
collection of 'pairs of socks' and not be able to select one sock from each pair. 
There can also be infinite sets where every subset is finite or has finite complement.

Another concept that is more standard is the concept of a 'hereditarily finite set'. 
For example, if $\mathbb N$ is the (infinite) set of natural numbers, the set $\{\mathbb N\}$, 
which has exactly one element (namely $\mathbb N$) is finite but it has an infinite set as an element, 
so it is not hereditarily finite. Similarly, the set $\{\{\mathbb N\}\}$, which has one element and that 
element has one element is finite but not hereditarily finite since the element of the element 
is an inifinite set.

\begin{remark}
More Reading materials
\begin{itemize}
    \item \href{https://math.stackexchange.com/questions/4925345/finite-infinite-set}{Finite and Infinite Sets}
    \item \href{https://math.stackexchange.com/questions/68659/set-theoretic-construction-of-the-natural-numbers}{Construction of Natural Numbers}
    \item \href{https://math.stackexchange.com/questions/3933535/what-is-the-definition-for-an-infinite-set}{Infinite Set}
\end{itemize}
\end{remark}





\bibliographystyle{alpha}
\bibliography{bib}
\end{document}
